% ==============================================================================
% Fichier : 03_recon_osint.tex
% Description : Reconnaissance Passive et OSINT
% ==============================================================================

\chapter{Reconnaissance Passive et OSINT}

La première règle du hacking est : "Connaître sa cible". Avant de lancer la moindre attaque technique, le hacker éthique collecte un maximum d'informations. Cette phase, appelée \textbf{Reconnaissance Passive}, a l'avantage d'être totalement invisible pour la cible (pas de paquets envoyés directement vers son réseau).

\section{Définition et Objectifs}
% ------------------------------------------------------------------------------

\subsection{OSINT (Open Source Intelligence)}

\begin{boxdef}
\textbf{OSINT :} Renseignement d'origine sources ouvertes. Il s'agit de collecter, analyser et exploiter des informations disponibles publiquement (Internet, médias, documents publics).
\end{boxdef}

\subsection{Objectifs de la Reconnaissance}

\begin{itemize}
    \item Identifier les \textbf{domaines} et sous-domaines appartenant à la cible.
    \item Découvrir les \textbf{adresses IP} publiques et les plages réseau.
    \item Identifier les \textbf{technologies} utilisées (Web Server, CMS, Cloud Provider).
    \item Trouver des \textbf{adresses email} et des \textbf{noms d'employés} (pour le Phishing ou le Brute Force).
    \item Découvrir des \textbf{fuites de données} (mots de passe passés, documents internes).
\end{itemize}

% ------------------------------------------------------------------------------
\section{Techniques et Outils de Collecte}
% ------------------------------------------------------------------------------

\subsection{Collecte d'Infrastructure}

\subsubsection{Whois et Registres}
Le protocole Whois permet d'interroger les bases de données des registraires de noms de domaine pour obtenir des informations sur le propriétaire.

\begin{lstlisting}[language=bash, caption=Requete Whois]
whois lizbiz.com
# Revele : Nom de l'administrateur, contacts techniques, 
# date de creation, et surtout le fournisseur d'acces (ISP).
\end{lstlisting}

\subsubsection{DNS Enumeration}
Le DNS (Domain Name System) est l'annuaire du web. Il contient des enregistrements précieux.

\begin{table}[h!]
\centering
\begin{tabular}{ll}
    \toprule
    \textbf{Type} & \textbf{Information révélée} \\
    \midrule
    A & Adresse IPv4 d'un hôte. \\
    MX & Serveurs de messagerie (Cibles privilégiées pour le phishing). \\
    TXT & Informations diverses (vérification propriété, SPF). \\
    NS & Serveurs de noms (DNS Primaire/Secondaire). \\
    \bottomrule
\end{tabular}
\caption{Enregistrements DNS critiques}
\end{table}

\begin{boxlizbiz}
\textbf{TP LiZbiZ :} Utilisez l'outil \texttt{dig} ou \texttt{nslookup} pour trouver l'adresse IP du serveur web de LiZbiZ et son serveur de messagerie.
\begin{lstlisting}[language=bash]
dig lizbiz.com ANY
\end{lstlisting}
\end{boxlizbiz}

\subsection{Moteurs de Recherche Spécialisés}

\subsubsection{Shodan : Le Moteur de Recherche des Objets Connectés}
Shodan scanne Internet en permanence et indexe les services (bannières) qu'il trouve.

\textbf{Recherche typique :}
\begin{itemize}
    \item \texttt{hostname:"lizbiz.com"} : Tous les services liés au domaine.
    \item \texttt{port:3389} : Recherche de bureaux à distance (RDP) exposés.
    \item \texttt{product:"Apache Tomcat"} : Recherche de versions vulnérables.
\end{itemize}

\subsubsection{Google Dorks (Google Hacking)}
Utilisation des opérateurs avancés de Google pour trouver des informations cachées.

\begin{table}[h!]
\centering
\small
\begin{tabular}{p{4cm}p{6cm}}
    \toprule
    \textbf{Opérateur} & \textbf{Exemple et Usage} \\
    \midrule
    \texttt{site:} & \texttt{site:lizbiz.com} (Limiter la recherche au site cible). \\
    \texttt{filetype:} & \texttt{filetype:pdf "confidentiel"} (Chercher des docs PDF sensibles). \\
    \texttt{inurl:} & \texttt{inurl:admin} (Chercher des pages d'administration). \\
    \texttt{intitle:} & \texttt{intitle:"index of /"} (Trouver des répertoires ouverts). \\
    \texttt{cache:} & Voir une version supprimée d'une page (Cache Google). \\
    \bottomrule
\end{tabular}
\caption{Opérateurs Google Dorks principaux}
\end{table}

\begin{boxlizbiz}
\textbf{Cas LiZbiZ :}
Un \texttt{Google Dork} révèle un fichier Excel indexé par erreur : \texttt{site:lizbiz.com filetype:xls "employés"}. Ce fichier contient la liste des emails internes, fournissant une base parfaite pour une campagne de Phishing ultérieure.
\end{boxlizbiz}

\subsection{Collecte d'Informations Humaines (Social Engineering Prep)}

\subsubsection{LinkedIn Scraping}
Les réseaux sociaux sont une mine d'or. LinkedIn révèle :
\begin{itemize}
    \item Les administrateurs système (cibles privilégiées).
    \item Les technologies utilisées (ex: "Administrateur Windows Server 2019 chez LiZbiZ").
    \item La structure hiérarchique.
\end{itemize}

\subsubsection{theHarvester}
Outil automatisé pour collecter emails, sous-domaines et noms d'hôtes à partir de multiples sources publiques.

\begin{lstlisting}[language=bash, caption=Utilisation de theHarvester]
theharvester -d lizbiz.com -l 500 -b google,linkedin
# -d : domaine
# -l : limite de resultats
# -b : source (moteur de recherche)
\end{lstlisting}

% ------------------------------------------------------------------------------
\section{Analyse de Métadonnées}
% ------------------------------------------------------------------------------

Les documents publics (PDF, Word, Images) publiés sur le site web contiennent souvent des métadonnées invisibles.

\subsection{ExifTool}
\begin{boxdef}
\textbf{Métadonnées :} Données cachées dans un fichier (Auteur, Logiciel utilisé, Date de création, Chemin du fichier source).
\end{boxdef}

\begin{boxlizbiz}
\textbf{Analyse LiZbiZ :}
Téléchargez le document "Rapport Annuel LiZbiZ.pdf". Utilisez \texttt{exiftool} pour l'analyser.
\begin{lstlisting}[language=bash]
exiftool rapport_lizbiz.pdf
# Resultat : "Creator Tool: Microsoft Word 2016", 
# "Author: Admin-IT". 
# Indice : Le compte "Admin-IT" existe probablement.
\end{lstlisting}
\end{boxlizbiz}